\Askhsh{35478}{4}
\begin{ylh}{1}
    \item 2.2 Παρουσίαση Στατιστικών Δεδομένων
    \item 2.3 Μέτρα Θέσης και Διασποράς
\end{ylh}
\begin{center}
\vspace{-1.5em}
\begin{tikzpicture}
\begin{axis}[
    axis lines=left,
    width=8cm, height=7cm,
    xmin=0, xmax=50,
    ymin=0, ymax=115,
    xtick={5, 15, 25, 35, 45},
    ytick={50},
    yticklabels={$50\%$},
    extra x ticks={0},
    extra x tick labels={$O$},
    ticklabel style={font=\small},
    ylabel={$F_i\%$},
    ylabel style={rotate=-90, at={(ticklabel cs:0.5)}, anchor=south, yshift=-1em},
    axis line style={-},
    clip=false
]
    % Bars using representative data for visual purposes
    \addplot[ybar interval, fill=maincolor!50, draw=black] coordinates {(5,25) (15,50) (25,80) (35,100) (45,0)};

    % Points on the top-left of each bar interval
    \addplot[only marks, mark=*, mark size=1.2pt] coordinates {(5,25) (15,50) (25,80) (35,100)};

    % Custom labels and ticks for other y-values
    \node[left] at (axis cs:0, 25) {$F_1\%$};
    \draw (axis cs:0,25) -- (axis cs:-0.5,25);
    \node[left] at (axis cs:0, 80) {$F_3\%$};
    \draw (axis cs:0,80) -- (axis cs:-0.5,80);
    \node[left] at (axis cs:0, 100) {$F_4\%$};
    \draw (axis cs:0,100) -- (axis cs:-0.5,100);
\end{axis}
\end{tikzpicture}
\end{center}
Οι χρόνοι σε λεπτά, που χρειάστηκαν οι μαθητές μιας τάξης για να λύσουν ένα μαθηματικό πρόβλημα ανήκουν στο διάστημα $[5,45)$ και έχουν ομαδοποιηθεί σε τέσσερις κλάσεις ίσου πλάτους. Τα δεδομένα των χρόνων εμφανίζονται στο παρακάτω ιστόγραμμα αθροιστικών σχετικών συχνοτήτων επί τοις εκατό.

\begin{alist}
\item Με βάση το παραπάνω ιστόγραμμα αθροιστικών σχετικών συχνοτήτων επί τοις εκατό, να υπολογίσετε τη διάμεσο των χρόνων που χρειάστηκαν οι μαθητές.
\monades{10}

\item
\begin{rlist}
	\item Στον επόμενο πίνακα συχνοτήτων της κατανομής των χρόνων, να αποδείξετε ότι $\alpha=8$.
	\monades{7}
	\item και να μεταφέρετε τον πίνακα κατάλληλα συμπληρωμένο στο τετράδιό σας.
\end{rlist}
\begin{center}
\begin{mytblr}{}
{Χρόνοι \\ (λεπτά)} & $x_i$ & $\nu_i$ & $f_i\%$ & $N_i$ & $F_i\%$ \\
$[5,15)$ & & $\alpha+4$ & & & \\
$[15,25)$ & & $3\alpha-6$ & & & \\
$[25,35)$ & & $2\alpha+8$ & & & \\
$[35,45)$ & & $\alpha-2$ & & & \\
Σύνολο & & & & & \\
\end{mytblr}
\end{center}\monades{8}
\end{alist}

